\documentclass[12pt]{article}
%\usepackage{Sweave}
\usepackage[margin=1in]{geometry}
%\usepackage[parfill]{parskip}
%\usepackage[hidelinks]{hyperref}
\usepackage[colorlinks = true,
            linkcolor = black,
            urlcolor  = blue,
            citecolor = black,
            anchorcolor = black]{hyperref}
\usepackage{rotating}
%\usepackage[backend=biber, style=apa, natbib=true]{biblatex}
\usepackage[citestyle=authoryear, style = apa, natbib=true, backend=bibtex, maxcitenames=2]{biblatex}
\addbibresource{references.bib}
\newtheorem{hyp}{H}
\newtheorem{hypo}{Hypothesis}
\usepackage[compact]{titlesec}
%\usepackage{titlesec}
\titleformat*{\section}{\bfseries\fontsize{14}{18}\selectfont}
\titleformat*{\subsection}{\bfseries\fontsize{12}{16}\selectfont}
%\usepackage{natbib}
\usepackage{csquotes}
\usepackage{xfrac}
\usepackage{mathtools}
\usepackage{booktabs}
\usepackage{lscape}
\usepackage{amsmath}
\usepackage[export]{adjustbox}
\usepackage{graphics}
\usepackage{array}
\usepackage[font=small,skip=1pt]{caption}
\usepackage{subcaption} 
\usepackage{ragged2e}
\usepackage{float}
\usepackage[section]{placeins}
\usepackage{hyperref}
\usepackage{multicol}
\usepackage{setspace}
\usepackage[hang,flushmargin]{footmisc}
\usepackage{afterpage}
\usepackage{soul}
\usepackage{xr}
\usepackage{booktabs}
\usepackage{multirow}
\usepackage{siunitx}
\usepackage{caption}
\usepackage{threeparttable}
\usepackage{makecell} 
\usepackage{tabularx}
\newcommand{\tabitem}{\textbullet~~}
\allowdisplaybreaks

% Set caption to be left-aligned
\captionsetup{justification=raggedright, singlelinecheck=false}

% Configure siunitx
\sisetup{
  group-separator = {,},
  group-minimum-digits = 4 % Start grouping from 4 digits
}

\doublespacing

\title{%
  The Signaling Gap: \\
  \large State-Owned Media and the Strategic Management of Sentiment during the Russian Invasion of Ukraine}
\author{Donald Moratz and Sarah Gerstein}
\date{\today}

%\title{The Signaling Gap: State-Owned Media and the Strategic Management of Sentiment during the Russian Invasion of Ukraine}

\begin{document}

\maketitle

%%%%%%%%%%%%%%%%%%%%%%%%%%%%%%% Second Rewrite %%%%%%%%%%%%%%%%%%%%%%%%%%%%%%%%%%%%%%%%%%

\abstract{Alliance commitments are designed to facilitate cooperation, but they can also prevent states from expressing disagreement through formal diplomatic channels, even when they genuinely dissent. We argue that state-owned media provides a mechanism for constrained states to signal genuine preferences to foreign governments, occupying a middle ground between formal channels that are too visible and private diplomacy that lacks credibility. Analyzing state media coverage before and after the 2022 invasion of Ukraine, we find that the invasion produced a larger negative shift in sentiment toward Russia in states with Russian security agreements than in neutral states. This occurred despite these same states refusing to formally condemn the invasion. Our analysis bridges a gap between international relations scholarship, which has largely treated state media as a domestic tool, and Western intelligence services, which have monitored it as a window into regime preferences since the Cold War.}


%%%%%%%%%%%%%%%%%%%%%%%%%%%%%%% Rewrite %%%%%%%%%%%%%%%%%%%%%%%%%%%%%%%%%%%%%%%%%%

%\abstract{Alliance commitments can prevent states from expressing disagreement through formal diplomatic channels, even when they genuinely dissent. We argue that state-owned media provides a mechanism for constrained states to signal genuine preferences to foreign governments, occupying a middle ground between formal channels that are too visible and private diplomacy that lacks credibility. Analyzing several hundred thousand articles from twelve states in Russia's near abroad before and after the 2022 invasion of Ukraine, we find that the invasion produced a larger negative shift in sentiment toward Russia in state-owned media of states with Russian security agreements than in neutral states. This occurred despite these same states refusing to formally condemn the invasion. Our analysis bridges a gap between international relations scholarship, which has largely treated state media as a domestic tool, and Western intelligence services, which have monitored it as a window into regime preferences since the Cold War.}


%%%%%%%%%%%%%%%%%%%%%%%%%%%%%%% Original %%%%%%%%%%%%%%%%%%%%%%%%%%%%%%%%%%%%%%%%%

%\abstract{How do states strategically signal positions during international crises? This paper investigates the 2022 Russian invasion of Ukraine as a seminal shock to analyze the divergent responses of state-controlled and independent media in post-Soviet and satellite states. While independent media reflects unconstrained sentiment, we argue that state-owned media serves as a strategic signaling device. Utilizing a corpus of several hundred thousand articles, we employ a Large Language Model (GPT-4o) to isolate sentiment toward Russia from general war-related topical shifts—an approach that outperforms traditional BERT-based classifiers in capturing nuanced elite cues. Our findings reveal that while sentiment declined across our sample, the drop was significantly muted in state-owned outlets. Crucially, we find that states bound by Russian security agreements utilized media sentiment to signal strategic dissatisfaction even when formal diplomatic condemnation remained politically impossible. These results highlight the role of controlled media in communicating "true" intent during geopolitical upheaval and demonstrate the power of generative AI in identifying latent elite signaling within massive text corpora.}

%%%%%%%%%%%%%%%%%%%%%%%%%%%%%% Third Rewrite %%%%%%%%%%%%%%%%%%%%%%%%%%%%

\section*{Introduction}

States that depend on a powerful patron for their security do not always have the luxury of saying what they think. Formal diplomatic channels---UN votes, official statements, public breaks---become liabilities when using them risks retaliation from the very power that guarantees a state's survival. Yet when a patron's actions draw international condemnation, silence risks being seen as endorsement. The Russian invasion of Ukraine on February 24, 2022---the largest land war in Europe since World War II---confronted Russia's formal security partners with exactly this bind. States bound to Russia through the Collective Security Treaty Organization and other formal security agreements had genuine reasons to oppose the invasion. It violated the sovereignty norms that protect small states and threatened their substantial economic ties with the West. But open condemnation of a security patron that had already demonstrated its willingness to use force in its near abroad was not a viable option. This is not a new problem. During the Cold War, Soviet-aligned states used state-owned media to communicate positions that formal diplomatic channels could not carry---and Western intelligence analysts learned to read those signals as indicators of genuine preferences \citep{cohen_theatre_1987}. We argue that the same channel Soviet-aligned states once used to signal under constraint remains operative today, and that the 2022 invasion reveals its logic with unusual clarity.
 
International relations scholarship has predominantly studied state media as a tool of domestic audience management \citep{mccombs_agenda-setting_2002, guriev_spin_2022}. Western governments, however, have monitored state media as a source of international intelligence for decades \citep{favere_debating_2024, bbc_monitoring_written_2016}, treating deviations from official positions as signals of genuine regime preferences \citep{williams_defining_2018}. We build on this practitioner insight to argue that state-owned media functions as an informal signaling channel through which constrained states communicate genuine preferences to external audiences---particularly Western governments. The credibility of this channel derives from the nature of state-owned media itself. External observers understand these outlets to be under government control, which means that when Kazakh or Belarusian state media shifts against Russia, foreign analysts do not interpret this as independent journalism but as a regime choice. This monitoring has only intensified in recent decades as open-source intelligence has become formalized within diplomatic and intelligence institutions \citep{manor_digitalization_2019, usic_ic_2024}. States in our sample are aware of this scrutiny, giving senders reason to believe their signals will be received and correctly interpreted. Furthermore, the signal is costly because allowing negative coverage of a security patron risks provoking retaliation from Moscow, a real cost that makes bluffing irrational. Crucially, state-owned media also offers what formal diplomatic channels cannot: deniability. The institutional separation between the government and state-owned media outlets gives states a tool to manage the relationship with Moscow if confronted.
 
Why state media rather than other channels? States facing this dilemma have a range of tools available, and the choice between them is not whether to signal, but a strategic calculus that must balance visibility and risk. The most public instruments---UN votes, formal diplomatic statements---are highly visible and directly attributable, which is precisely what makes them too costly for states bound by security agreements. The least visible alternative---quiet, private diplomatic meetings---avoids provocation but sacrifices credibility; a message delivered behind closed doors is unverifiable by the audiences it is meant to reach and easily dismissed as cheap talk \citep{fearon_signaling_1997}. State-owned media occupies the middle ground. It is public enough to be observed by foreign governments that systematically monitor it, yet it preserves enough institutional distance from official policy that the sending state retains plausible deniability.
 
We test this argument by exploiting variation in formal security relationships with Russia across twelve states in Russia's near abroad, drawing on a corpus of several hundred thousand newspaper articles published before and after the invasion.\footnote{The twelve states span three categories of alignment. Russia-aligned: Armenia, Azerbaijan, Belarus, and Kazakhstan. Neutral: Georgia, Moldova, Kosovo, and Serbia. NATO-aligned: Albania, Hungary, and Turkey. Ukraine is included in the dataset but excluded from the main analysis given its direct role as the invaded party.} Four states were bound to Russia through the Collective Security Treaty Organization or equivalent bilateral security agreements at the time of the invasion. Four occupied broadly neutral positions, maintaining distance from both Russian security structures and full Western integration. The remaining states are NATO members. We employ a generative AI model (GPT-4o) to measure sentiment toward Russia at the article level, validating performance against a human-coded gold standard. Our identification strategy uses the sharp discontinuity of February 24 to estimate the causal effect of the invasion on media sentiment within each group through a regression discontinuity in time design. To compare across groups, we leverage a Regression Difference-in-Discontinuity Design (RDID). The sharp discontinuity of the invasion substantially reduces confounding concerns, though comparisons across groups remain correlational in nature. We subject these comparisons to extensive robustness checks, which we detail in our methods section.
 
Our findings build sequentially. We first establish that the invasion produced a genuine negative shock to media sentiment toward Russia across the region---without this baseline, the divergence we identify in subsequent tests would be uninterpretable. We then show that this decline was not uniform: independent media responded significantly more negatively than state-owned media, consistent with strategic management of coverage by state outlets. The critical test of our argument comes in the comparison across security relationships. A straightforward domestic audience management account would predict that Russia-aligned states suppress negative coverage to protect the alliance relationship. Instead, we find the opposite: state-owned media in states with Russian security agreements shifted more negatively toward Russia than state-owned media in neutral states, despite these same states refusing to condemn the invasion through formal diplomatic channels. This pattern is consistent with constrained states using media to signal genuine dissatisfaction to external audiences precisely because formal channels were unavailable.
 
This paper makes two primary contributions. First, our analysis demonstrates that the signals practitioners have long monitored are theoretically structured and empirically identifiable. The patterns we document are not idiosyncratic to a single crisis but follow directly from the logic of alliance constraints: states whose formal channels are constrained by dependence on a powerful patron produce systematically different media signals than states whose channels remain open. This suggests that state media sentiment is not merely noise but a legible instrument of international communication, one that the discipline can study with the same rigor it applies to formal diplomatic behavior. Second, we identify an understudied mechanism of informal signaling under alliance constraints. The signaling literature has focused on how states communicate resolve to adversaries or commitment to allies through formal channels; we demonstrate that when those channels are foreclosed, states turn to informal ones, and that the resulting signals follow predictable patterns rooted in the structure of their security commitments. Methodologically, our approach pairs generative AI-based sentiment measurement with a regression discontinuity design, offering a template for large-scale analysis of media signals in contexts where hand-coding is infeasible.
 
The article proceeds as follows. We first develop our theoretical argument and derive three hypotheses about media divergence under alliance constraints. We then present the case, justifying our sample and identification strategy. The following sections describe our data and research design before presenting results. We conclude by discussing implications for how states signal preferences under geopolitical constraint.

%%%%%%%%%%%%%%%%%%%%%%%%%%%%%% Second Rewrite %%%%%%%%%%%%%%%%%%%%%%%%%%%%

%\section*{Introduction}

%States that depend on a powerful patron for their security do not always have the luxury of saying what they think. Formal diplomatic channels---UN votes, official statements, public breaks---become liabilities when using them risks retaliation from the very power that guarantees a state's survival. Yet when a patron's actions draw international condemnation, silence risks being seen as endorsement. The Russian invasion of Ukraine on February 24, 2022---the largest land war in Europe since World War II---confronted Russia's formal security partners with exactly this bind. States bound to Russia through the Collective Security Treaty Organization and other formal security agreements had genuine reasons to oppose the invasion. It violated the sovereignty norms that protect small states and threatened their substantial economic ties with the West. But open condemnation of a security patron that had already demonstrated its willingness to use force in its near abroad was not a viable option. This is not a new problem. During the Cold War, Soviet-aligned states used state-owned media to communicate positions that formal diplomatic channels could not carry---and Western intelligence analysts learned to read those signals as indicators of genuine preferences \citep{cohen_theatre_1987}. We argue that the same channel Soviet-aligned states once used to signal under constraint remains operative today, and that the 2022 invasion reveals its logic with unusual clarity.


%States that depend on a powerful patron for their security do not always have the luxury of saying what they think. Formal diplomatic channels---UN votes, official statements, public breaks---become liabilities when using them risks retaliation from the very power that guarantees a state's survival. Yet inaction risks being seen as endorsement. The Russian invasion of Ukraine on February 24, 2022 confronted Russia's formal security partners with exactly this bind. States bound to Russia through the Collective Security Treaty Organization and other formal security agreements had genuine reasons to oppose the largest land war in Europe since World War II---the invasion violated the sovereignty norms that protect small states and threatened their substantial economic ties with the West. But open condemnation of a security patron that had already demonstrated its willingness to use force in its near abroad was not a viable option. This is not a new problem. During the Cold War, Soviet-aligned states used state-owned media to communicate positions that formal diplomatic channels could not carry---and Western intelligence analysts learned to read those signals as indicators of genuine preferences \citep{cohen_theatre_1987}. We argue that the same channel Soviet-aligned states once used to signal under constraint remains operative today, and that the 2022 invasion reveals its logic with unusual clarity.
 
%We contend that state-owned media functions as an informal signaling channel through which constrained states communicate genuine preferences to external audiences---particularly Western governments. The credibility of this channel derives from the nature of state-owned media itself. External observers understand these outlets to be under government control, which means that deviations from official positions are informative; when Kazakh or Belarusian state media shifts against Russia, foreign analysts do not interpret this as independent journalism but as a regime choice. Western governments have monitored Soviet and post-Soviet state media for decades \citep{favere_debating_2024, williams_defining_2018}, a practice that has only intensified as open-source intelligence has become formalized within diplomatic and intelligence institutions \citep{manor_digitalization_2019, usic_ic_2024}. This monitoring is common knowledge among the states in our sample, giving senders reason to believe their signals will be received and correctly interpreted. Furthermore, the signal is costly because allowing negative coverage of a security patron risks provoking retaliation from Moscow, a real cost that makes bluffing irrational. Yet state-owned media offers one additional advantage over formal diplomatic channels: deniability. The institutional separation between the government and state-owned media outlets gives states a tool to manage the relationship with Moscow if confronted.

%Why state media rather than other channels? States facing this dilemma have a range of tools available, and the choice between them is not whether to signal, but a strategic calculus that must balance visibility and risk. The most public instruments---UN votes, formal diplomatic statements---are highly visible and directly attributable, which is precisely what makes them too costly for states bound by security agreements. The least visible alternative---quiet, private diplomatic meetings---avoids provocation but sacrifices credibility; a message delivered behind closed doors is unverifiable by the audiences it is meant to reach and easily dismissed as cheap talk \citep{fearon_signaling_1997}. State-owned media occupies the middle ground. It is public enough to be observed by foreign governments that systematically monitor it, yet it preserves enough institutional distance from official policy that the sending state retains plausible deniability.
 
%Why do states use state controlled media rather than other channels? The choice is not whether to signal, but a strategic calculus that must balance visibility and risk. The most public instruments---UN votes, formal diplomatic statements---are highly visible and directly attributable, which is precisely what makes them too costly for states bound by security agreements. The least visible alternative---quiet, private diplomatic meetings---avoids provocation but sacrifices credibility; a message delivered behind closed doors is unverifiable by the audiences it is meant to reach and easily dismissed as cheap talk \citep{fearon_signaling_1997}. State-owned media occupies the middle ground. It is public enough to be observed by foreign governments that systematically monitor it, yet it preserves enough institutional distance from official policy that the sending state retains plausible deniability.
 
%We test this argument by exploiting variation in formal security relationships with Russia across twelve states in Russia's near abroad, drawing on a corpus of several hundred thousand newspaper articles published before and after the invasion.\footnote{The twelve states span three categories of alignment. Russia-aligned: Armenia, Azerbaijan, Belarus, and Kazakhstan. Neutral: Georgia, Moldova, Kosovo, and Serbia. NATO-aligned: Albania, Hungary, and Turkey. Ukraine is included in the dataset but excluded from the main analysis given its direct role as the invaded party.} Four states were bound to Russia through the Collective Security Treaty Organization or equivalent bilateral security agreements at the time of the invasion. Four occupied broadly neutral positions, maintaining distance from both Russian security structures and full Western integration. The remaining states are NATO members. We employ a generative AI model (GPT-4o) to measure sentiment toward Russia at the article level, validating performance against a human-coded gold standard. Our identification strategy uses the sharp discontinuity of February 24 to estimate the causal effect of the invasion on media sentiment within each group through a regression discontinuity in time design. To compare across groups, we leverage a Regression Difference-in-Discontinuity Design (RDID). The sharp discontinuity of the invasion substantially reduces confounding concerns, though comparisons across groups remain correlational in nature. We subject these comparisons to extensive robustness checks, which we detail in our methods section.
 
%Our findings build sequentially. We first establish that the invasion produced a genuine negative shock to media sentiment toward Russia across the region---without this baseline, the divergence we identify in subsequent tests would be uninterpretable. We then show that this decline was not uniform: independent media responded significantly more negatively than state-owned media, consistent with strategic management of coverage by state outlets. The critical test of our argument comes in the comparison across security relationships. A straightforward domestic audience management account would predict that Russia-aligned states suppress negative coverage to protect the alliance relationship. Instead, we find the opposite: state-owned media in states with Russian security agreements shifted more negatively toward Russia than state-owned media in neutral states, despite these same states refusing to condemn the invasion through formal diplomatic channels. This pattern is consistent with constrained states using media to signal genuine dissatisfaction to external audiences precisely because formal channels were unavailable.
 
%This paper contributes to international relations scholarship in two primary ways. First, we bridge a gap between international relations theory and practice. Intelligence analysts and foreign governments devote substantial resources to monitoring state-owned media as a window into regime preferences---a practice with institutional roots stretching back to the Cold War---yet the discipline has predominantly studied state media as a tool of domestic audience management. We show that state media also operates as an instrument of international communication, and that the signals it carries are theoretically interpretable and empirically measurable. Second, we identify an understudied mechanism of informal signaling under alliance constraints. The signaling literature has focused on how states communicate resolve to adversaries or commitment to allies through formal channels; we demonstrate that when those channels are foreclosed, states turn to informal ones, and that the resulting signals follow predictable patterns rooted in the structure of their security commitments. Methodologically, our approach pairs generative AI-based sentiment measurement with a regression discontinuity design, offering a template for large-scale analysis of media signals in contexts where hand-coding is infeasible.
 
%The article proceeds as follows. We first develop our theoretical argument and derive three hypotheses about media divergence under alliance constraints. We then present the case, justifying our sample and identification strategy. The following sections describe our data and research design before presenting results. We conclude by discussing implications for how states signal preferences under geopolitical constraint.

%%%%%%%%%%%%%%%%%%%%%%%%%%%%%% Rewrite %%%%%%%%%%%%%%%%%%%%%%%%%%%%%%%%%%%

%\subsection*{Rewrite}

%How do constrained states signal genuine preferences when formal diplomatic channels are too costly to use? The Russian invasion of Ukraine in February 2022 serves as a seminal shock to the post-Cold War era. As NATO debated invasion preparations \citep{jakes_for_2024} and US political rhetoric questioned alliance commitments \citep{sanger_outburst_2024}, smaller states faced acute pressure to signal their strategic alignments without triggering costly retaliation from Moscow. One observable mechanism through which such signals can be transmitted is state media, which shapes public opinion through selective attribution and sentiment \citep{guriev_theory_2020, rozenas_how_2019}, yet the international relations literature has focused primarily on how these elite cues impact domestic audiences \citep{guisinger_mapping_2017}. Consequently, the divergence between state and private media sentiment, and the specific utility of media as an international signaling device, remains under-explored. We argue that state-owned media functions as a deniable informal signaling channel for states whose formal diplomatic constraints make official condemnation impossible, and that this dynamic is observable in the sentiment divergence between state and independent media during the 2022 invasion.

%By using a generative AI model rather than traditional BERT-based classifiers, we isolate sentiment specifically toward Russia while controlling for topic. This ensures we capture shifts in attitude toward Russia itself, separate from the confounding influence of general war reporting. This approach allows us to identify how strategic elites responded to the crisis and clarifies the relationship between elites and foreign policy. We present three findings. First, the invasion caused a significant decline in media sentiment across all topics, separate from the shift toward war coverage. Second, we find that the causal effect of the invasion on sentiment differs significantly between independent and state-owned media, highlighting the strategic management of sentiment. Finally, we show that states with Russian security agreements expressed greater negative sentiment post-invasion than neutral countries, despite refusing to condemn the invasion publicly. This suggests these states utilize media to communicate genuine preferences when formal statements are politically impossible.

%Signaling is a cornerstone of international relations literature. While theoretical work often centers on signaling credibility to rivals \citep{fearon_signaling_1997} or allies \citep{blankenship_trivial_2022}, research has largely focused on the domestic agenda-setting role of media \citep{mccombs_agenda-setting_2002}. Building on \citet{cohen_theatre_1987}, we contend that states strategically control sentiment toward Russia not only for domestic audiences but also to signal preferences to the international community. Such media signals reveal genuine preferences when formal diplomatic channels, such as UN votes, are constrained by high political costs. The invasion of Ukraine represents a crisis where states may publicly maintain one position while signaling a different intent through controlled media. This strategic behavior stands in contrast to independent media which remains unconstrained by state signaling imperatives.

%Previously, data limitations hindered large-scale examination of these dynamics. However, advances in natural language processing provide new avenues for inquiry. We study how state-controlled media in states in Russia's near abroad responded to the invasion compared to independent sources. Utilizing a corpus of several hundred thousand articles from the period surrounding February 24, 2022, we applied a keyword-based approach to identify Russia-centric coverage. We then employed a generative AI model (GPT-4o) to conduct sentiment analysis, validating the model's performance against a human-coded gold standard. This approach allows us to causally identify how coverage changed around the discontinuity of the invasion for both state-owned and independent media, while controlling for topical shifts. The comparison across states with and without Russian security agreements is correlational, though the sharp discontinuity of the invasion provides a natural experiment that substantially reduces concerns about confounding within each group. %Our validation yielded a Krippendorff's alpha of 0.86, indicating high inter-coder reliability on an ordinal sentiment scale. 

%The article proceeds as follows. We first develop our theoretical argument and derive three hypotheses about media divergence under alliance constraints. We then present the case, justifying our sample and identification strategy. The following sections describe our data and research design before presenting results. We conclude by discussing implications for how states signal preferences under geopolitical constraint.

%%%%%%%%%%%%%%%%%%%%%%%%%%%%%% Original %%%%%%%%%%%%%%%%%%%%%%%%%%%%%%%%%%

%\subsection*{Original}

%How do media sources respond to significant disruptions in the international status quo? The Russian invasion of Ukraine in February 2022 serves as a seminal shock to the post-Cold War era. As NATO debated invasion preparations \citep{jakes_for_2024} and US political rhetoric questioned alliance commitments \citep{sanger_outburst_2024}, it is critical to study how states signaled their evolving strategic alignments. While state media is recognized as a primary tool for influencing public opinion \citep{guriev_theory_2020} through selective attribution and sentiment \citep{rozenas_how_2019}, the international relations literature has focused primarily on how these elite cues impact domestic audiences \citep{guisinger_mapping_2017}. Consequently, the divergence between state and private media sentiment, and the specific utility of media as an international signaling device, remains under-explored.

%Signaling is a cornerstone of international relations literature. While theoretical work often centers on signaling credibility to rivals \citep{fearon_signaling_1997} or allies \citep{blankenship_trivial_2022}, research has largely focused on the domestic agenda-setting role of media \citep{mccombs_agenda-setting_2002}. In this article, we argue that the role of state-owned media as an international signaling device has been overlooked. Building on \citet{cohen_theatre_1987}, we contend that states strategically control sentiment toward Russia not only for domestic audiences but also to signal preferences to the international community. Such media signals reveal true intentions when formal diplomatic channels, such as UN votes, are constrained by high political costs. The invasion of Ukraine represents a crisis where states may publicly maintain one position while signaling a different intent through controlled media. This strategic behavior stands in contrast to independent media which remains unconstrained by state signaling imperatives.

%Previously, data limitations hindered large-scale examination of these dynamics. However, advances in natural language processing provide new avenues for inquiry. We study how state-controlled media in former Soviet, satellite, and allied states responded to the invasion compared to independent sources. Utilizing a corpus of several hundred thousand articles from the period surrounding February 24, 2022, we applied a keyword-based approach to identify Russia-centric coverage. We then employed a generative AI model (GPT-4o) to conduct sentiment analysis, validating the model's performance against a human-coded gold standard. Our validation yielded a Krippendorff's alpha of 0.86, indicating high inter-coder reliability on an ordinal sentiment scale. This approach allows us to causally identify how coverage changed around the discontinuity of the invasion while controlling for topical shifts.

%By using a generative AI model rather than traditional BERT-based classifiers, we isolate sentiment specifically toward Russia while controlling for topic. This ensures we capture shifts in attitude toward Russia itself, separate from the confounding influence of general war reporting. This approach allows us to identify how strategic elites responded to the crisis and clarifies the relationship between elites and foreign policy. We present three findings. First, the invasion caused a significant decline in media sentiment across all topics, separate from the shift toward war coverage. Second, we find a differential causal effect between independent and state-owned media, highlighting the strategic management of sentiment. Finally, we show that states with Russian security agreements expressed greater negative sentiment post-invasion than neutral countries, despite refusing to condemn the invasion publicly. This suggests these states utilize media to communicate true feelings when formal statements are politically impossible.

\newpage

\printbibliography

%\bibliographystyle{apalike}
%\bibliography{references.bib}

\end{document}